\documentclass[a4paper,11pt]{article}

% ─── Geometry ────────────────────────────────────────────────────────────────
\usepackage{geometry}
\geometry{margin=1in, headheight=15pt}

% ─── Mathematics ─────────────────────────────────────────────────────────────
\usepackage{amsmath}
\usepackage{amssymb}
\usepackage{mathtools}

% ─── Graphics / Colours ──────────────────────────────────────────────────────
\usepackage{graphicx}
\usepackage{xcolor}
\definecolor{blue3}{RGB}{0,0,180}
\definecolor{green3}{RGB}{0,120,0}
\definecolor{orange3}{RGB}{200,100,0}
\definecolor{red3}{RGB}{180,0,0}
\definecolor{grey3}{RGB}{80,80,80}

% ─── Units ───────────────────────────────────────────────────────────────────
\usepackage{siunitx}
\sisetup{detect-all}

% ─── Links ───────────────────────────────────────────────────────────────────
\usepackage[colorlinks=true,linkcolor=blue3,citecolor=blue3,urlcolor=blue3]{hyperref}
\usepackage{cleveref}

% ─── Boxes ───────────────────────────────────────────────────────────────────
\usepackage{tcolorbox}
\tcbuselibrary{skins,breakable}

\newtcolorbox{answerbox}[1][]{
    colback=red3!5,
    colframe=red3!60,
    coltitle=red3,
    title={Solution:},
    fonttitle=\bfseries,
    breakable,
    #1
}

% ─── Tables / Lists ──────────────────────────────────────────────────────────
\usepackage{booktabs}
\usepackage{enumitem}

% ─── Notation commands ───────────────────────────────────────────────────────
\newcommand{\ktilde}{\tilde{k}}
\newcommand{\ytilde}{\tilde{y}}
\newcommand{\ctilde}{\tilde{c}}
\newcommand{\htilde}{\tilde{h}}
\newcommand{\dd}{\mathrm{d}}

% ─── Header / Footer ─────────────────────────────────────────────────────────
\usepackage{fancyhdr}
\pagestyle{fancy}
\fancyhf{}
\lhead{\textbf{14E022 Macroeconomics I} \quad Practice Exam 2 -- Solutions}
\rhead{BSE 2025--26}
\cfoot{\thepage}

% ═════════════════════════════════════════════════════════════════════════════
\begin{document}
% ═════════════════════════════════════════════════════════════════════════════

\begin{center}
    {\LARGE \textbf{14E022 Macroeconomics I}}\\[6pt]
    {\Large \textbf{Practice Exam 2 --- Answer Key}}\\[8pt]
    {\large BSE --- Academic Year 2025--26}
\end{center}

\noindent\rule{\textwidth}{0.4pt}

% ═════════════════════════════════════════════════════════════════════════════
\section{Question 1 --- Augmented Solow Model with Human Capital}
% ═════════════════════════════════════════════════════════════════════════════

\subsection*{1.\ (5 points)}

\begin{answerbox}
\textbf{Output per effective worker.}

\[
    \ytilde = \frac{Y}{AL}
    = \frac{K^{\alpha}H^{\beta}(AL)^{1-\alpha-\beta}}{AL}
    = \frac{K^{\alpha}}{(AL)^{\alpha}}\cdot\frac{H^{\beta}}{(AL)^{\beta}}
    = \ktilde^{\alpha}\htilde^{\beta}. \quad\checkmark
\]

\textbf{Law of motion for $\ktilde$.}

Since $\ktilde = K/(AL)$, differentiate by the quotient/product rule:
\[
    \dot{\ktilde} = \frac{\dot{K}}{AL} - K\,\frac{\dot{A}L + A\dot{L}}{(AL)^2}
    = \frac{\dot{K}}{AL} - \ktilde\!\left(\frac{\dot{A}}{A} + \frac{\dot{L}}{L}\right)
    = \frac{\dot{K}}{AL} - (g + n)\,\ktilde.
\]
Substituting $\dot{K} = s_K Y - \delta K$:
\[
    \dot{\ktilde}
    = \frac{s_K Y - \delta K}{AL}
    - (n+g)\,\ktilde
    = s_K\ytilde - \delta\ktilde - (n+g)\ktilde.
\]
\[
    \boxed{\dot{\ktilde} = s_K\,\ktilde^{\alpha}\htilde^{\beta}
    - (n+g+\delta)\,\ktilde.}
\]

\textbf{Law of motion for $\htilde$.}  An identical calculation with $\dot{H} = s_H Y - \delta H$ gives:
\[
    \boxed{\dot{\htilde} = s_H\,\ktilde^{\alpha}\htilde^{\beta}
    - (n+g+\delta)\,\htilde.}
\]
\end{answerbox}

\subsection*{2.\ (7 points)}

\begin{answerbox}
\textbf{(a)} Setting $\dot{\ktilde}=0$ and $\dot{\htilde}=0$:
\begin{align*}
    s_K\,\ktilde^{*\alpha}\htilde^{*\beta} &= (n+g+\delta)\,\ktilde^*, \\
    s_H\,\ktilde^{*\alpha}\htilde^{*\beta} &= (n+g+\delta)\,\htilde^*.
\end{align*}
Dividing the first by the second:
\[
    \frac{s_K}{s_H} = \frac{\ktilde^*}{\htilde^*}
    \quad\Longrightarrow\quad
    \boxed{\htilde^* = \frac{s_H}{s_K}\,\ktilde^*.}
\]

\textbf{(b)} Substituting $\htilde^* = (s_H/s_K)\ktilde^*$ into the first steady-state equation:
\[
    s_K\,\ktilde^{*\alpha}\left(\frac{s_H}{s_K}\right)^{\!\beta}\ktilde^{*\beta}
    = (n+g+\delta)\,\ktilde^*.
\]
Collect powers of $s_K$ and $s_H$ and of $\ktilde^*$:
\[
    s_K\cdot s_K^{-\beta}\cdot s_H^{\beta}\cdot \ktilde^{*(\alpha+\beta)}
    = (n+g+\delta)\,\ktilde^*.
\]
\[
    s_K^{1-\beta}\,s_H^{\beta}\,\ktilde^{*(\alpha+\beta-1)} = (n+g+\delta).
\]
Raise both sides to the power $1/(\alpha+\beta-1)$.  Since $\alpha+\beta < 1$, the exponent
$\alpha+\beta-1 < 0$, so $1/(\alpha+\beta-1) = -1/(1-\alpha-\beta)$:
\[
    \ktilde^* = \left[\frac{s_K^{1-\beta}\,s_H^{\beta}}{n+g+\delta}
    \right]^{1/(1-\alpha-\beta)}.
\]
Using $\htilde^* = (s_H/s_K)\ktilde^*$:
\[
    \htilde^* = \frac{s_H}{s_K}\left[\frac{s_K^{1-\beta}\,s_H^{\beta}}{n+g+\delta}
    \right]^{1/(1-\alpha-\beta)}
    = \left[\frac{s_K^{\alpha}\,s_H^{1-\alpha}}{n+g+\delta}
    \right]^{1/(1-\alpha-\beta)}.
\]
\[
    \boxed{\ktilde^* = \left[\frac{s_K^{1-\beta}\,s_H^{\beta}}{n+g+\delta}
    \right]^{1/(1-\alpha-\beta)},\qquad
    \htilde^* = \left[\frac{s_K^{\alpha}\,s_H^{1-\alpha}}{n+g+\delta}
    \right]^{1/(1-\alpha-\beta)}.}
\]

Steady-state output per effective worker:
\[
    \ytilde^* = \ktilde^{*\alpha}\htilde^{*\beta}
    = \left[\frac{s_K^{1-\beta}\,s_H^{\beta}}{n+g+\delta}\right]^{\!\alpha/(1-\alpha-\beta)}
      \cdot\left[\frac{s_K^{\alpha}\,s_H^{1-\alpha}}{n+g+\delta}\right]^{\!\beta/(1-\alpha-\beta)}.
\]
Combining exponents ($s_K$ exponent: $(1-\beta)\alpha/D + \alpha\beta/D = \alpha/D$;
$s_H$ exponent: $\beta\alpha/D + (1-\alpha)\beta/D = \beta/D$, where $D = 1-\alpha-\beta$):
\[
    \boxed{\ytilde^* = \left[\frac{s_K^{\alpha}\,s_H^{\beta}}{(n+g+\delta)^{\alpha+\beta}}
    \right]^{1/(1-\alpha-\beta)}.}
\]
\end{answerbox}

\subsection*{3.\ (7 points)}

\begin{answerbox}
Since $y^* = A_t\,\ytilde^*$:
\[
    \ln y^* = \ln A_0 + g\,t
    + \frac{\alpha}{1-\alpha-\beta}\ln s_K
    + \frac{\beta}{1-\alpha-\beta}\ln s_H
    - \frac{\alpha+\beta}{1-\alpha-\beta}\ln(n+g+\delta).
\]

\textbf{(a)} Elasticity of $y^*$ w.r.t.\ $s_K$:
\begin{itemize}
    \item \textbf{Basic Solow} ($\beta=0$): $\varepsilon = \alpha/(1-\alpha)$.
    For $\alpha=1/3$: $\varepsilon = 1/2$.
    \item \textbf{Augmented Solow} ($\alpha=\beta=1/3$):
    $\varepsilon = (1/3)/(1-1/3-1/3) = (1/3)/(1/3) = 1$.
\end{itemize}
The elasticity \textbf{doubles} from $0.5$ to $1.0$.

\textbf{(b)} The augmented model amplifies the effect of $s_K$ differences (because
$\alpha+\beta > \alpha$ reduces the denominator $1-\alpha-\beta$) and adds a
\emph{second channel} through $s_H$.  Cross-country variation in both physical and
human capital investment rates can now explain a much larger fraction of income
differences.  MRW showed this model can explain roughly 80\% of cross-country income
variation (vs.\ $\sim$60\% for basic Solow), substantially reducing the unexplained
residual attributed to TFP.
\end{answerbox}

\subsection*{4.\ (6 points)}

\begin{answerbox}
\textbf{Key limitation:} The model treats $s_H$ (human capital investment rate) as
exogenous.  In reality, $s_H$ is itself an endogenous choice driven by institutions,
returns to education, credit constraints, and government policy.  Countries with low
$s_H$ may have low $s_H$ \emph{because} of their low income or weak institutions, not
as an independent causal factor.  This creates an \textbf{endogeneity problem}: the
model can ``explain'' income differences with $s_H$, but $s_H$ itself needs to be
explained.  A deeper theory would endogenise the education decision (e.g.\ the
Uzawa-Lucas model) to understand \emph{why} $s_H$ differs across countries.

Additionally, TFP ($A$) is treated as equal across countries.  Residual TFP differences
(even after adding $H$) remain large and may reflect institutional quality,
misallocation, or technology adoption barriers --- factors the augmented Solow model
does not address.
\end{answerbox}

\clearpage
% ═════════════════════════════════════════════════════════════════════════════
\section{Question 2 --- Ramsey Model: Permanent Increase in TFP Growth}
% ═════════════════════════════════════════════════════════════════════════════

\subsection*{1.\ (5 points)}

\begin{answerbox}
\textbf{(a)} Law of motion:
\[
    \dot{\ktilde} = \ktilde^{\alpha} - \ctilde - (n + x + \delta)\,\ktilde.
\]

\textbf{(b)} Euler equation (per-effective-worker):
\[
    \frac{\dot{\ctilde}}{\ctilde}
    = \frac{1}{\sigma}\bigl(\alpha\ktilde^{\alpha-1} - \delta - \rho - \sigma x\bigr).
\]

\textbf{(c)} Loci:
\begin{itemize}
    \item $\dot{\ctilde}=0$: vertical line at $\ktilde^*$ where
    $\alpha\ktilde^{*(\alpha-1)} = \delta + \rho + \sigma x$.
    To the left, $\dot{\ctilde}>0$; to the right, $\dot{\ctilde}<0$.
    \item $\dot{\ktilde}=0$: $\ctilde = \ktilde^{\alpha} - (n+x+\delta)\ktilde$
    (hump-shaped curve).  Above it, $\dot{\ktilde}<0$; below, $\dot{\ktilde}>0$.
\end{itemize}
The steady state is at the intersection.  The saddle path approaches it from both sides.
\end{answerbox}

\subsection*{2.\ (7 points)}

\begin{answerbox}
\textbf{(a)} Modified Golden Rule: $\alpha\ktilde^{*(\alpha-1)} = \delta + \rho + \sigma x$.
\[
    \boxed{\ktilde^* = \left[\frac{\alpha}{\delta + \rho + \sigma x}
    \right]^{1/(1-\alpha)}.}
\]
\[
    \ctilde^* = \ktilde^{*\alpha} - (n+x+\delta)\,\ktilde^*.
\]

\textbf{(b)} Modified Golden Rule: $f'(\ktilde^*) = \delta + \rho + \sigma x$.

\textbf{(c)} The TVC (in per-effective-worker form) is
$\lim_{t\to\infty}\lambda(t)\,\ktilde(t)\,e^{-\rho'\!t} = 0$,
where $\rho' = \rho - n - (1-\sigma)x$.

On the BGP, $\ktilde$ is constant (growth rate 0), and from the FOC $\lambda = \ctilde^{-\sigma}$
with $\ctilde$ growing at rate $\dot{\ctilde}/\ctilde = (1/\sigma)(\alpha\ktilde^{*(\alpha-1)}-\delta-\rho-\sigma x)$.
On the BGP this equals 0 (since we are at the steady state $\ktilde^*$), so $\lambda$ is
constant.  Thus $\lambda(t)\ktilde(t) \to \text{const}$, and the TVC requires
$e^{-\rho' t}\to 0$, i.e.\ $\rho' > 0$:
\[
    \rho' = \rho - n - (1-\sigma)x > 0
    \quad\Longleftrightarrow\quad
    \rho - n > (1-\sigma)x.
\]
This ensures that the discount factor dominates growth in the costate and that lifetime
utility is bounded (otherwise the household's objective integral diverges).
\end{answerbox}

\subsection*{3.\ (7 points)}

\begin{answerbox}
\textbf{(a)} New steady state with $x' > x$:
\[
    \ktilde^{*\prime} = \left[\frac{\alpha}{\delta + \rho + \sigma x'}
    \right]^{1/(1-\alpha)} < \ktilde^*
    = \left[\frac{\alpha}{\delta + \rho + \sigma x}\right]^{1/(1-\alpha)}
\]
since $\delta + \rho + \sigma x' > \delta + \rho + \sigma x$ and the exponent
$1/(1-\alpha) > 0$, so the expression decreases.

For per-effective-worker consumption, use the $\dot{\ktilde}=0$ locus evaluated at the
steady state:
\[
    \ctilde^{*\prime} = \ktilde^{*\prime\alpha} - (n+x'+\delta)\,\ktilde^{*\prime}
    = \ktilde^{*\prime}\bigl[\ktilde^{*\prime\alpha-1} - (n+x'+\delta)\bigr].
\]
Comparing with $\ctilde^* = \ktilde^*\bigl[\ktilde^{*\alpha-1} - (n+x+\delta)\bigr]$,
there are two reinforcing effects both reducing $\ctilde^{*\prime}$:
\begin{enumerate}[label=(\roman*), itemsep=2pt]
    \item \textbf{Lower $\ktilde^{*\prime}$}: since $\ktilde^{*\prime} < \ktilde^*$,
    the production term $\ktilde^{*\prime\alpha} < \ktilde^{*\alpha}$.
    \item \textbf{Higher effective depreciation}: $(n+x'+\delta) > (n+x+\delta)$,
    so more output is required to keep $\ktilde$ constant.
\end{enumerate}
Both channels reduce $\ctilde^{*\prime}$ below $\ctilde^*$.  Indeed, since $\ktilde^{*\prime}$
satisfies the Modified Golden Rule $\alpha\ktilde^{*\prime\alpha-1} = \delta+\rho+\sigma x'$,
we can also write $\ctilde^{*\prime} = \ktilde^{*\prime\alpha}(1 - (n+x'+\delta)/
\ktilde^{*\prime\alpha-1}\cdot\alpha^{-1}\cdot\alpha)$; the key point is that $\ktilde^{*\prime}$
is strictly smaller, and the $\dot{\ktilde}=0$ locus is shifted down, confirming
$\ctilde^{*\prime} < \ctilde^*$.

\textbf{Intuition:} Faster TFP growth means that per-effective-worker capital depreciates
faster (the ``effective worker'' grows faster), so the economy sustains a lower level
of $\ktilde$ in steady state.  The household also requires a higher return to save (to
compensate for faster-growing consumption), which further discourages capital accumulation
in effective units.

\textbf{(b) Shifts:}
\begin{itemize}
    \item $\dot{\ctilde}=0$ shifts \textbf{left}: from $\ktilde^*$ to
    $\ktilde^{*\prime} < \ktilde^*$.
    \item $\dot{\ktilde}=0$ shifts \textbf{down}: higher $(n+x'+\delta)$ steepens the
    effective depreciation line, lowering the curve
    $\ctilde = \ktilde^{\alpha} - (n+x'+\delta)\ktilde$.
\end{itemize}

\textbf{(c)}
\begin{itemize}
    \item $\ktilde$ does \textbf{not jump}: $K$, $A$, and $L$ are all continuous at $t=0$
    (only $\dot{A}/A$ changes, not $A$ itself).
    \item $\ctilde$ \textbf{jumps up}: the household immediately raises per-effective-worker
    consumption to reach the new saddle path.  Since $\ktilde(0) = \ktilde^* > \ktilde^{*\prime}$,
    the economy is to the right of the new steady state.
    \item \textbf{Transition}: both $\ctilde$ and $\ktilde$ fall gradually along the new
    saddle path from $(\ktilde^*, \ctilde(0^+))$ to $(\ktilde^{*\prime}, \ctilde^{*\prime})$.
\end{itemize}
\end{answerbox}

\subsection*{4.\ (6 points)}

\begin{answerbox}
Per-capita consumption and income:
\[
    c_t = \ctilde_t\,A_t, \qquad y_t = \ytilde_t\,A_t.
\]

Before the shock, $c_t$ and $y_t$ grow at rate $x$.  After the shock, they grow at rate
$x' > x$ in the long run (since $\ctilde_t \to \ctilde^{*\prime}$ and $A_t$ grows at $x'$).

Although $\ctilde^{*\prime} < \ctilde^*$ (the per-effective-worker \emph{level} falls),
the per-capita \emph{growth rate} is permanently higher.  At some finite time, the new
path overtakes the old one:
\[
    \ctilde^{*\prime}\,A_0 e^{x' t} > \ctilde^*\,A_0 e^{x t}
    \quad\text{for all } t > \frac{\ln(\ctilde^*/\ctilde^{*\prime})}{x'-x}.
\]

The household is better off because the present discounted value of lifetime utility is
higher: faster permanent growth outweighs the short-run reduction in consumption level.
Faster technological progress is unambiguously welfare-improving.
\end{answerbox}

\clearpage
% ═════════════════════════════════════════════════════════════════════════════
\section{Question 3 --- Development Accounting and Misallocation}
% ═════════════════════════════════════════════════════════════════════════════

\subsection*{1.\ (5 points)}

\begin{answerbox}
\textbf{(a)} Take logs: $\ln Y = \ln A + \alpha\ln K + (1-\alpha)\ln L$.

Differentiate w.r.t.\ time:
\[
    \boxed{\gamma_Y = \gamma_A + \alpha\,\gamma_K + (1-\alpha)\,\gamma_L.}
\]

The \textbf{Solow residual} is:
$\gamma_A = \gamma_Y - \alpha\gamma_K - (1-\alpha)\gamma_L$.
It captures everything that contributes to output growth beyond measured factor
accumulation --- including technology, efficiency, measurement error, and institutional
change.

\textbf{(b)} Two factors besides technology: (i) \textbf{misallocation} of resources
across firms/sectors (changes in allocative efficiency show up as TFP changes);
(ii) \textbf{mismeasurement} of input quality --- e.g.\ not adjusting labour for
human capital or capital for quality improvements.
\end{answerbox}

\subsection*{2.\ (7 points)}

\begin{answerbox}
In per-worker terms: $y = A k^{\alpha}$.

Taking ratios:
$\ln(y_R/y_P) = \ln(A_R/A_P) + \alpha\ln(k_R/k_P)$.

With $y_R/y_P = 27$, $k_R/k_P = 27$, $\alpha = 1/3$:

\textbf{(a)} Capital contribution:
$\alpha\ln(k_R/k_P) = \tfrac{1}{3}\ln 27 = \tfrac{1}{3}\cdot 3\ln 3 = \ln 3$.

As a fraction of $\ln(y_R/y_P) = \ln 27 = 3\ln 3$:
\[
    \text{Capital share} = \frac{\ln 3}{3\ln 3} = \boxed{\frac{1}{3} \approx 33\%.}
\]

\textbf{(b)} TFP contribution (residual):
$\ln(A_R/A_P) = 3\ln 3 - \ln 3 = 2\ln 3$.
\[
    \text{TFP share} = \frac{2\ln 3}{3\ln 3} = \boxed{\frac{2}{3} \approx 67\%.}
\]

\textbf{(c)} TFP differences are the \textbf{dominant} proximate source of cross-country
income differences, accounting for about two-thirds of the gap.  This motivates studying
what drives TFP: technology adoption, institutions, and misallocation.
\end{answerbox}

\subsection*{3.\ (7 points)}

\begin{answerbox}
\textbf{(a)} Efficient allocation equalises MPK across firms:
$\alpha A k_1^{\alpha-1} = \alpha A k_2^{\alpha-1}$ $\Rightarrow$ $k_1 = k_2 = K/2$.
\[
    Y^{\mathrm{eff}} = 2A(K/2)^{\alpha} = 2^{1-\alpha}A\,K^{\alpha}.
\]

\textbf{(b)} With the subsidy, firm~1 optimises taking its effective rental cost as
$(1-\tau)r$; firm~2 takes $r$.  Profit-maximising capital demands equate MPK to
effective rental cost:
\begin{align*}
    \text{Firm 1:}\quad \alpha A k_1^{\alpha-1} &= (1-\tau)r, \\
    \text{Firm 2:}\quad \alpha A k_2^{\alpha-1} &= r.
\end{align*}
Divide the first equation by the second ($\alpha A$ cancels):
\[
    \frac{k_1^{\alpha-1}}{k_2^{\alpha-1}} = (1-\tau)
    \quad\Longrightarrow\quad
    \left(\frac{k_1}{k_2}\right)^{\!\alpha-1} = (1-\tau).
\]
Raise both sides to the power $\tfrac{1}{\alpha-1}$.  Since $\alpha < 1$,
$\alpha - 1 < 0$, so $\tfrac{1}{\alpha-1} = -\tfrac{1}{1-\alpha}$:
\[
    \frac{k_1}{k_2}
    = (1-\tau)^{1/(\alpha-1)}
    = (1-\tau)^{-1/(1-\alpha)}.
\]
Because $\tau \in (0,1)$, we have $(1-\tau) \in (0,1)$, so $(1-\tau)^{-1/(1-\alpha)} > 1$.
\[
    \boxed{\frac{k_1}{k_2} = (1-\tau)^{-1/(1-\alpha)} > 1.}
\]
Firm 1 (subsidised) receives \emph{too much} capital relative to the efficient allocation.

\textbf{(c)} Aggregate output with misallocation:
$Y^{\mathrm{mis}} = A(k_1^{\alpha} + k_2^{\alpha})$ with $k_1 \neq k_2$ but
$k_1 + k_2 = K$.

Since $f(k) = k^{\alpha}$ is \textbf{strictly concave} (for $\alpha < 1$), Jensen's
inequality gives:
\[
    \frac{k_1^{\alpha} + k_2^{\alpha}}{2} < \left(\frac{k_1 + k_2}{2}\right)^{\!\alpha}
    = (K/2)^{\alpha}
\]
whenever $k_1 \neq k_2$.  Therefore:
\[
    Y^{\mathrm{mis}} = A(k_1^{\alpha} + k_2^{\alpha}) < 2A(K/2)^{\alpha} = Y^{\mathrm{eff}}.
    \quad\checkmark
\]

\textbf{Intuition:} The subsidy channels too much capital to firm~1, where the marginal
product is low, and too little to firm~2, where the marginal product is high.  Aggregate
output would be higher if capital were reallocated from firm~1 to firm~2 until marginal
products are equalised.  This is the \textbf{essence of misallocation}: dispersion in
marginal products signals that a reallocation of inputs could increase aggregate output
without using more resources.
\end{answerbox}

\subsection*{4.\ (6 points)}

\begin{answerbox}
\textbf{(a)} \textbf{Misallocation} is defined as the dispersion of marginal revenue
products across firms: $\text{Var}(\text{MRPK}_i) > 0$.  In an efficient economy with no
frictions, all firms equalise MRPK to the common rental rate $r$, so
$\text{Var}(\text{MRPK}) = 0$.  Positive variance means that some firms have higher
returns to capital than others, implying that total output could be increased by
reallocating capital from low-MRPK firms to high-MRPK firms.

\textbf{(b)} Two concrete examples of policies/frictions creating capital wedges:
\begin{enumerate}[label=(\roman*)]
    \item \textbf{Directed/subsidised credit:} State-owned banks in China provide
    below-market loans to state-owned enterprises (SOEs), giving them artificially low
    capital costs.  Private firms face higher borrowing rates $\Rightarrow$ MRPK dispersion.
    \item \textbf{Size-dependent regulations:} In India, labour laws (e.g.\ the Industrial
    Disputes Act) impose high firing costs on firms above a size threshold.  This
    discourages productive firms from growing to efficient scale, keeping capital trapped
    in smaller, less productive establishments.
\end{enumerate}

Other valid examples: trade protection for specific sectors, entry barriers from
licensing, corruption/connections determining access to resources, land-use regulations
preventing reallocation.
\end{answerbox}

\clearpage
% ═════════════════════════════════════════════════════════════════════════════
\section{Short Questions}
% ═════════════════════════════════════════════════════════════════════════════

\subsection*{1.\ (9 points) --- Golden Rule vs.\ Modified Golden Rule}

\begin{answerbox}
\textbf{(a)} In the Solow model, steady-state consumption is
$c^* = f(\ktilde^*) - (n+g+\delta)\ktilde^*$.  Maximising over $\ktilde^*$:
\[
    \frac{\dd c^*}{\dd\ktilde^*} = f'(\ktilde^*_{GR}) - (n+g+\delta) = 0
    \quad\Longrightarrow\quad
    \boxed{f'(\ktilde^*_{GR}) = n + g + \delta.}
\]
This is \emph{independent of $s$} because $s$ determines which steady state the economy
reaches, but the Golden Rule picks the $\ktilde^*$ that maximises $c^*$ regardless of
how you get there.

\textbf{(b)} In the Ramsey model, the Modified Golden Rule is:
$f'(\ktilde^*) = \delta + \rho + \sigma x$.

This generally differs from the Golden Rule because the Ramsey household is
\emph{impatient} ($\rho > 0$): it discounts future consumption, so it does not save
enough to reach the Golden Rule level.  The two coincide if and only if
$\rho + \sigma x = n + g$, i.e.\ $\rho = n + (1-\sigma)g$ (with $g=x$), which is
a knife-edge condition.

\textbf{(c)} The Solow economy \emph{can} be dynamically inefficient: if $s$ is high
enough that $\ktilde^* > \ktilde^*_{GR}$, the economy over-accumulates capital (MPK $< n+g+\delta$).
Reducing $s$ would raise consumption in every period.

The Ramsey economy \emph{cannot} be dynamically inefficient: optimising households never
accumulate beyond the Modified Golden Rule, which satisfies $f'(\ktilde^*) > n+g+\delta$
(given the TVC constraint $\rho - n > (1-\sigma)x$).  Dynamic efficiency is guaranteed
by forward-looking optimisation.
\end{answerbox}

\subsection*{2.\ (8 points) --- Role of Monopolistic Competition}

\begin{answerbox}
\textbf{(a)} Ideas are \textbf{non-rival}: once created, they can be used by everyone at
zero marginal cost.  Under perfect competition, price equals marginal cost, so $p = MC = 0$
for ideas.  But then innovators earn zero revenue and cannot recoup the fixed cost of R\&D.
No firm would invest in innovation $\Rightarrow$ no endogenous growth.  The fundamental
problem: non-rivalry requires increasing returns, which is incompatible with perfect
competition.

\textbf{(b)} Monopolistic competition gives innovators \emph{temporary market power}.
In the Romer model, the markup $p = 1/\alpha > 1$ (above marginal cost of production)
generates positive profits $\pi > 0$.  These profits are the \emph{reward for innovation}
and are essential for financing R\&D.  Without the markup, the value of creating a new
variety $V = \pi/r = 0$, and no one would innovate.

\textbf{(c)} The markup creates a \textbf{static inefficiency}: monopolists restrict
output below the socially optimal level ($x^{\text{mkt}} < x^{\text{plnr}}$).  A
\textbf{subsidy on intermediate purchases} equal to $1 - \alpha$ (so the effective price
faced by final good producers falls from $1/\alpha$ to $1$) corrects this distortion.
This raises intermediate usage to the efficient level while preserving the innovator's
profit incentive (profits are replaced by the subsidy).
\end{answerbox}

\subsection*{3.\ (8 points) --- Expanding Varieties vs.\ Quality Ladders}

\begin{answerbox}
\textbf{(a)} In Romer, new varieties \emph{coexist} with old ones --- no firm is displaced.
In Aghion--Howitt, new quality improvements \emph{replace} old producers.  Key implication:
in quality ladders, incumbent monopolists \textbf{lose their rents} when a successor
innovates.  This creates a direct conflict between incumbents and entrants that is absent
in the expanding varieties framework.

\textbf{(b)} \textbf{Creative destruction}: the process by which new innovators displace
incumbents.  The \textbf{business-stealing externality}: the entrant captures the entire
market, destroying the incumbent's profit stream, but does not internalise this cost.
This leads to \textbf{over-innovation} relative to the social optimum, because the private
return to innovation includes ``stealing'' rents that were already generating social value.

\textbf{(c)} The \textbf{appropriability externality}: the innovator cannot capture all
the social value of the innovation --- consumers benefit from quality improvements beyond
what monopoly profits reflect (consumer surplus).  This pushes toward
\textbf{under-innovation}.

The two externalities act in \emph{opposite} directions.  The net effect is ambiguous and
depends on parameters.  This means the market innovation rate can be either too high or
too low relative to the social optimum.  Optimal policy requires balancing both
externalities, potentially with multiple instruments (R\&D subsidies/taxes + output subsidies).
\end{answerbox}

\vspace{12pt}
\noindent\rule{\textwidth}{0.4pt}
\begin{center}\textit{End of Answer Key --- Practice Exam 2.}\end{center}

\end{document}
